% problem_id: S001-lemma-cohomology-de-rham-base-change
% source_id: S001 (Stacks Project)
% source_file: perfect.tex
% statement_locator: perfect.tex lines 5553--5574
% proof_locator: perfect.tex lines 5576--5678
% env: lemma
% pairing: tight (verified)
% status: complete (statement + author proof verbatim + reason + locators)
%
\begin{lemma}
\label{lemma-cohomology-de-rham-base-change}
Let $a : X \to S$ be a quasi-compact and quasi-separated morphism
of schemes. Let $\mathcal{F}^\bullet$ be a locally bounded
complex of $a^{-1}\mathcal{O}_S$-modules. Assume for all $n \in \mathbf{Z}$
the sheaf $\mathcal{F}^n$ is a flat $a^{-1}\mathcal{O}_S$-module and
$\mathcal{F}^n$ has the structure of a quasi-coherent $\mathcal{O}_X$-module
compatible with the given $a^{-1}\mathcal{O}_S$-module structure (but the
differentials in the complex $\mathcal{F}^\bullet$ need not
be $\mathcal{O}_X$-linear). Then the following hold
\begin{enumerate}
\item $Ra_*\mathcal{F}^\bullet$ is locally bounded,
\item $Ra_*\mathcal{F}^\bullet$ is in $D_\QCoh(\mathcal{O}_S)$,
\item $Ra_*\mathcal{F}^\bullet$ locally has finite tor dimension,
\item $\mathcal{G} \otimes_{\mathcal{O}_S}^\mathbf{L} Ra_*\mathcal{F}^\bullet =
Ra_*(a^{-1}\mathcal{G} \otimes_{a^{-1}\mathcal{O}_S} \mathcal{F}^\bullet)$
for $\mathcal{G} \in \QCoh(\mathcal{O}_S)$, and
\item $K \otimes_{\mathcal{O}_S}^\mathbf{L} Ra_*\mathcal{F}^\bullet =
Ra_*(a^{-1}K \otimes_{a^{-1}\mathcal{O}_S}^\mathbf{L} \mathcal{F}^\bullet)$
for $K \in D_\QCoh(\mathcal{O}_S)$.
\end{enumerate}
\end{lemma}

\begin{proof}
Parts (1), (2), (3) are local on $S$ hence we may and do assume $S$
is affine. Since $a$ is quasi-compact we conclude that $X$ is quasi-compact.
Since $\mathcal{F}^\bullet$ is locally bounded, we conclude that
$\mathcal{F}^\bullet$ is bounded.

\medskip\noindent
For (1) and (2) we can use the first spectral sequence
$R^pa_*\mathcal{F}^q \Rightarrow R^{p + q}a_*\mathcal{F}^\bullet$ of
Derived Categories, Lemma \ref{derived-lemma-two-ss-complex-functor}.
Combining Cohomology of Schemes, Lemma
\ref{coherent-lemma-quasi-coherence-higher-direct-images}
and Homology, Lemma \ref{homology-lemma-biregular-ss-converges}
we conclude.

\medskip\noindent
Let us prove (3) by the induction principle of
Cohomology of Schemes, Lemma \ref{coherent-lemma-induction-principle}.
Namely, for a quasi-compact open of $U$ of $X$ consider the
condition that $R(a|_U)_*(\mathcal{F}^\bullet|_U)$ has
finite tor dimension. If $U, V$ are quasi-compact open in
$X$, then we have a relative Mayer-Vietoris distinguished triangle
$$
R(a|_{U \cup V})_*\mathcal{F}^\bullet|_{U \cup V} \to
R(a|_U)_*\mathcal{F}^\bullet|_U \oplus
R(a|_V)_*\mathcal{F}^\bullet|_V \to
R(a|_{U \cap V})_*\mathcal{F}^\bullet|_{U \cap V} \to
$$
by Cohomology, Lemma \ref{cohomology-lemma-unbounded-relative-mayer-vietoris}.
By the behaviour of tor amplitude in distinguished triangles
(see Cohomology, Lemma \ref{cohomology-lemma-cone-tor-amplitude})
we see that if we know the result for $U$, $V$, $U \cap V$, then
the result holds for $U \cup V$. This reduces us to the case where
$X$ is affine. In this case we have
$$
Ra_*\mathcal{F}^\bullet = a_*\mathcal{F}^\bullet
$$
by Leray's acyclicity lemma
(Derived Categories, Lemma \ref{derived-lemma-leray-acyclicity})
and the vanishing of higher direct images of quasi-coherent modules
under an affine morphism
(Cohomology of Schemes, Lemma \ref{coherent-lemma-relative-affine-vanishing}).
Since $\mathcal{F}^n$ is $S$-flat by assumption and $X$ affine, the modules
$a_*\mathcal{F}^n$ are flat for all $n$. Hence $a_*\mathcal{F}^\bullet$
is a bounded complex of flat $\mathcal{O}_S$-modules and hence has
finite tor dimension.

\medskip\noindent
Proof of part (5). Denote
$a' : (X, a^{-1}\mathcal{O}_S) \to (S, \mathcal{O}_S)$
the obvious flat morphism of ringed spaces. Part (5) says that
$$
K \otimes_{\mathcal{O}_S}^\mathbf{L} Ra'_*\mathcal{F}^\bullet =
Ra'_*(L(a')^*K \otimes_{a^{-1}\mathcal{O}_S}^\mathbf{L}
\mathcal{F}^\bullet)
$$
Thus
Cohomology, Equation (\ref{cohomology-equation-projection-formula-map})
gives a functorial map from the left to the right and we want to show
this map is an isomorphism.
This question is local on $S$ hence we may and do assume $S$
is affine. The rest of the proof is {\it exactly} the same as the
proof of Lemma \ref{lemma-cohomology-base-change} except that we have
to show that the functor
$K \mapsto Ra'_*(L(a')^*K \otimes_{a^{-1}\mathcal{O}_S}^\mathbf{L}
\mathcal{F}^\bullet)$ commutes with direct sums.
This is where we will use $\mathcal{F}^n$ has the structure
of a quasi-coherent $\mathcal{O}_X$-module. Namely, observe that
$K \mapsto L(a')^*K
\otimes_{a^{-1}\mathcal{O}_S}^\mathbf{L} \mathcal{F}^\bullet$
commutes with arbitrary direct sums. Next, if
$\mathcal{F}^\bullet$ consists of a single quasi-coherent
$\mathcal{O}_X$-module $\mathcal{F}^\bullet = \mathcal{F}^n[-n]$
then we have $L(a')^*G
\otimes_{a^{-1}\mathcal{O}_S}^\mathbf{L} \mathcal{F}^\bullet =
La^*K \otimes_{\mathcal{O}_X}^\mathbf{L} \mathcal{F}^n[-n]$, see
Cohomology, Lemma \ref{cohomology-lemma-variant-derived-pullback}.
Hence in this case the commutation with direct sums follows from
Lemma \ref{lemma-quasi-coherence-pushforward-direct-sums}.
Now, in general, since $S$ is affine (hence $X$ quasi-compact)
and $\mathcal{F}^\bullet$ is locally bounded, we see that
$$
\mathcal{F}^\bullet = (\mathcal{F}^a \to \ldots \to \mathcal{F}^b)
$$
is bounded. Arguing by induction on $b - a$ and considering the
distinguished triangle
$$
\mathcal{F}^b[-b] \to (\mathcal{F}^a \to \ldots \to \mathcal{F}^b)
\to (\mathcal{F}^a \to \ldots \to \mathcal{F}^{b - 1}) \to
\mathcal{F}^b[-b + 1]
$$
the proof of this part is finished. Some details omitted.

\medskip\noindent
Proof of part (4). Let $a' : (X, a^{-1}\mathcal{O}_S) \to (S, \mathcal{O}_S)$
be as above. Since $\mathcal{F}^\bullet$ is a locally bounded
complex of flat $a^{-1}\mathcal{O}_S$-modules we see the complex
$a^{-1}\mathcal{G} \otimes_{a^{-1}\mathcal{O}_S} \mathcal{F}^\bullet$
represents $L(a')^*\mathcal{G}
\otimes_{a^{-1}\mathcal{O}_S}^\mathbf{L}
\mathcal{F}^\bullet$ in $D(a^{-1}\mathcal{O}_S)$. Hence (4)
follows from (5).
\end{proof}
